\scanpage{57}

1) Il contient un él. d'ordre $4$, i.e. n'est pas formé d'él d'ordre $2$ et $3$.

Il \unsure{n'|ne} \deleted{est pas formé de} \inserted{peut contenir} deux él. d'ordre $2$, sinon $H$ serait un groupe diédral, ce qui n'est pas le cas. Donc il contiendrait deux él. d'ordre $3$, \deleted{\unsure{soit}} \inserted{qui engendrent $H$}

\deleted{$\rho_f,\rho_{\deleted{f'}g}$ $(f,g\in F)$ (\unsure{comme|jamais} $\rho,\rho'$ sont égaux ni opposés) On ne peut avoir \deleted{\unsure{$f=g'$}} $g=f$ ou $g=f'$, car $\rho_f,\rho_g$ \unsure{commuteraient}. Si on écrit $H\simeq\unsure{\mathfrak S_4|\mathcal S_4}$) $\pi_{123}$, et $\pi_{234}$ ou $\pi_{324}$ (modulo conjugaison) mais alors $\rho\rho'$ est}

\inserted{\unsure{or|ou}} les él. d'ordre $3$ sont dans $H^{+}$, absurde

2) Il contient aussi un él. d'ordre $3$, i.e. n'est pas formé \unsure{seult|seulmt|seulement} d'él. d'ordre $4$ et $2$. Car dans ce cas, comme il ne contient $2$ él. d'ordre $2$, il contiendrait au moins un autre él. d'ordre \deleted{\unsure{3}}$4$.

\deleted{Soient $\rho,\rho'$} \inserted{\unsure{SPS|Sps|OPS}} dans $(\rho,\rho',\rho'')$ que $\rho,\rho'$ sont d'ordre $4$ (quitte à utiliser $\mathfrak S_3$), donc de la forme $\rho_s,\rho_{\deleted{s'}t}$. \deleted{Mais \unsure{si}} \inserted{Évidt} \deleted{\unsure{a)}} $t\ne s,s'$ (antipodique de $s$), donc $(t,s)$ est un arc, on est dans le cas traité en c), où on a un él. d'ordre $3$.

3) Il reste (comme le cas de $2$ él. d'ordre $4$ est traité) à regarder le cas où on a un seul. él. d'ordre $4$, un d'ordre $3$, soit $(\rho,\rho',\rho'')$ avec

% archive sheet number: 56
